\documentclass[10pt,a4paper]{amsart}
\usepackage[margin=19mm]{geometry}
\usepackage{amsmath,amssymb,mathtools}
\usepackage[scr=rsfs,cal=euler]{mathalfa}
\usepackage{xfrac}
\usepackage[unicode,hidelinks,bookmarks=false]{hyperref}
\newcommand{\ie}{i.e.\ }
\newcommand{\cf}{cf.\ }
\newcommand{\resp}{respectively\ }
\newcommand{\Subsection}{\subsection}
\providecommand{\nobreakdash}{\nobreak\hbox{-}}
\setlength{\emergencystretch}{2em}
\multlinegap0pt
\usepackage{amssymb}
\usepackage{algorithm}\usepackage{algpseudocode}
\usepackage{float}
\newtheorem{lemma}{Lemma}
\title{The Stability of a Coupled Degenerate Wave System Under Boundary Control}
\author{SUN Ya-nan, ZHANG Qiong}
\date{Source: arXiv:2604.05441v1 (CC BY 4.0)}
\begin{document}
\maketitle

\noindent\emph{Source and licence.} The Stability of a Coupled Degenerate Wave System Under Boundary Control. Version arXiv:2604.05441v1 (CC BY 4.0), distributed by arXiv under the Creative Commons Attribution 4.0 International licence, http://creativecommons.org/licenses/by/4.0/, as recorded in the arXiv licence metadata for this exact version and shown on the abstract page https://arxiv.org/abs/2604.05441v1. Full licence notice: \texttt{licenses/arxiv-2604.05441-CC-BY-4.0.txt}.\par\smallskip

\noindent\textit{Editorial scope.} Counted unit: the article's lemma lemma (source0.tex lines 417--419). The article's complete original proof of it is delivered with it (source0.tex lines 420--528, one \texttt{proof} environment of 109 lines). No mathematics is written, repaired or re-derived: the statement and the proof are the article's own lines, in the article's own notation, and no retained line is substituted or re-flowed.  Retained uncounted as the source-local context the proof needs: lines 390--392.  Nothing was omitted from any retained span, so the package carries no omission record.  No retained line contains a citation, so the wrapper carries no bibliography and no bibliography entry.  No retained line contains a figure environment, an image-inclusion command or drawing code, so no image is delivered and the package declares no asset dependency.  Editorial formatting only, in addition: an AMS \texttt{amsart} preamble that restates the article's own declarations and macros used by the retained lines, namely the environments \texttt{lemma} on separate counters, together with the standard packages those lines need.  The retained environments print in this document as Lemma 1 in order of appearance, whereas the article prints the same items with the numbering of its own full document.  The complete native source of the article as distributed in the arXiv e-print for this exact version is bundled unchanged as \texttt{sources/197957/source0.tex}.
\begin{equation}\label{def}
\mathcal A_i U = (v,(au')',y,(bw')'),\;\;\forall\;U = (u,v,w,y),
\end{equation}

\begin{lemma}
The operator $\mathcal A_i$ generates a $C_0$-semigroup of contractions $T_i(t)$ on the energy space $\mathcal H_i$, and $0\in\rho(\mathcal A_i)$, the resolvent of the operator $\mathcal A_i$, $i=1,2$.
\end{lemma}

\begin{proof}
For every $U=(u,v,w,y)\in\mathcal D(\mathcal A_i)$, $i=1,2$ it follows from \eqref{def} that
$$\Re \langle \mathcal A_iU,U \rangle = - a(1)\left|v(1)\right|^2 \leq 0,$$
which implies that $\mathcal A_i$ is dissipative.

Next, we shall show that $0\in\rho(\mathcal A_i)$, $i=1,2$. Let $U\in\mathcal D(\mathcal A_1)$ satisfy $\mathcal A_1 U = 0$.  Then,
\begin{equation}
	\begin{cases}
		v=0, &x\in[0,1],\\
		(au')' = 0,&x\in[0,1],\\
		y = 0,&x\in[-1,0],\\
		(bw')' =0,&x\in[-1,0].
	\end{cases}
\end{equation}
Then, it follows from the transmission condition that
\begin{equation}\label{0-du}
	u' = (bw')(0)x^{-\mu_a},\;\;\forall\;x\in(0,1),
\end{equation}
and
\begin{equation}\label{0-dw}
	w' = (bw')(0)(1+x)^{-\mu_b},\;\;\forall\;x\in(-1,0).
\end{equation}
Since $w\in W_b^1(-1,0)$, it holds that $w(-1) = 0$. Then, for any $x\in(-1,0)$,
$$ w(x)= {(bw')(0)\over1-\mu_b}(1+x)^{1-\mu_b}. $$
By the transmission condition, one has, for any $x\in(0,1)$,
\begin{equation}
	\begin{split}
		u(x) =(bw')(0)\left({x^{1-\mu_a}\over 1-\mu_a} + {1\over1-\mu_b}\right).
	\end{split}
\end{equation}
Substituting the above results into the boundary condition, one has
$$\gamma u(1)+u'(1)+v(1) = (bw')(0) \left(1 + {\gamma\over1-\mu_a} +{\gamma\over1-\mu_b}\right) = 0,$$
which implies that $(bw')(0) = 0$, i.e., $U\equiv0$.

Let $U\in\mathcal D(\mathcal A_2)$ satisfy $\mathcal A_2U=0$, we have that \eqref{0-du} and \eqref{0-dw} remain valid. Since $1\leq\mu_b<2$, one has
\begin{equation*}
	\lim_{x\rightarrow -1^+} (bw')(x) = 0,
\end{equation*}
combining with the transmission condition, which implies that
$$au'\equiv 0,\;\forall\; x\in(0,1),$$and
$$bw'\equiv 0,\;\forall\; x\in(-1,0).$$
Then, it follows from $a(x)>0$ $(x\neq0)$, $b(x)>0$ $(x\neq-1)$ and the transmission condition that
$$u(x)\equiv u(1),\;\;\forall\;x\in(0,1),$$
and $$ w(x) \equiv u(1),\;\;\forall\; x\in(-1,0).$$
Substituting the above results into the boundary condition yields
$\gamma u(1) = 0$,
which implies that $U\equiv 0$.

Then, we prove that $\mathcal A_i\;(i=1,2)$ is surjective. For any fixed $F= (f_1,f_2,g_1,g_2)\in \mathcal H_1\setminus\{0\}$. Let $\mathcal A_1U = F$, i.e.,
\begin{subequations}
	\begin{align}
		&v=f_1,&&x\in[0,1],\label{0-a}\\
		&(au')'=f_2,&&x\in[0,1],\label{0-b}\\
		&y=g_1,&&x\in[-1,0],\label{0-c}\\
		&(bw')'= g_2,&&x\in[-1,0].\label{0-d}
	\end{align}
\end{subequations}
From \eqref{0-b} and the transmission condition, one has
\begin{equation}\label{0-u'}
	u'(x)={(bw')(0)\over a(x)}+{1\over a(x)}\int_0^x f_2(r)dr,\;\;\forall\;x\in(0,1),
\end{equation}
and
\begin{equation}\label{0-u}
	u(x) = u(0) + \int_0^x {1\over a(s)}\left((bw')(0)+\int_0^s f_2(r)dr\right) ds,\;\;\forall\;x\in(0,1).
\end{equation}
Since $w\in W_b^2(-1,0)$, it holds that $w(-1) =0$. Then,
\begin{equation}\label{0-w'}
	w'(x) = {(bw')(0)\over b(x)} - {1\over b(x)}\int_x^0 g_2(r)dr,\;\;\forall\;x\in(-1,0),
\end{equation}
and
\begin{equation}\label{0-w}
	w(x) = \int_{-1}^x {1\over b(s)} \left( (bw')(0)  -  \int_s^0 g_2(r)dr\right) ds,\;\;\forall\;x\in(-1,0).
\end{equation}
By the transmission conditions, one has that
\begin{equation}\label{0-u00}
	u(0) = w(0) = {(bw')(0)\over 1-\mu_b} - \int_{-1}^0{1\over b(s)} \int_s^0 g_2(r)dr ds.
\end{equation}
Substituting \eqref{0-a}, \eqref{0-u'} and \eqref{0-u} into the boundary conditions $\gamma u(1) + u'(1) +v(1) =0$ , one has that
\begin{equation}\label{0-au'00}
	\begin{split}
		(bw')(0) = &\left(1 + {\gamma\over1-\mu_b}+{\gamma\over1-\mu_a}\right)^{-1}
		\left(\gamma\int_{-1}^0 {1\over b(s)}\int_s^0 g_2(r) dr ds - f_1(1)\right.\\
		&\left.\quad-\int_0^1 f_2(r)dr
		- \gamma\int_0^1{1\over a(s)}\int_0^sf_2(r)dr ds \right).
	\end{split}
\end{equation}
Thus, $u(x)$ and $w(x)$ can be uniquely determined by \eqref{0-u}, \eqref{0-w}, \eqref{0-u00} and \eqref{0-au'00}.

Let $\mathcal A_2 U= F$, then \eqref{0-a}-\eqref{0-d}, \eqref{0-u'} and \eqref{0-u} remain valid. It follows from $\lim_{x\rightarrow -1^+} bw'= 0$ that
\begin{equation}
	(bw')(x) = \int_{-1}^x g_2(r)dr,\;\;\forall\;x\in(-1,0),
\end{equation}
and
\begin{equation}\label{0-w1}
	w(x) = w(0) - \int_x^0{1\over b(s)}\int_{-1}^s g_2(r) drds,\;\;\forall\;x\in(-1,0).
\end{equation}
By the transmission condition, one has that
\begin{equation}\label{0-au'0}
	\lim_{x\to0^+}(au')(x) = (bw')(0) = \int_{-1}^0 g_2(r)dr.
\end{equation}
Substituting \eqref{0-au'0} into \eqref{0-u'}, \eqref{0-u}, respectively, and combining with the boundary condition, we have that
\begin{equation}\label{0-u0}
	\begin{split}
		u(0) = -{1\over\gamma} \int_0^1 f_2(r)dr - \int_0^1 {1\over a(s)}\int_0^s f_2(r)dr ds
		- {1\over\gamma} f_1(1)-\left({1\over\gamma} + {1\over1-\mu_a}\right)\int_{-1}^0 g_2(r)dr.
	\end{split}
\end{equation}
In the end, $u(x)$ and $w(x)$ can be  uniquely determined by \eqref{0-u}, \eqref{0-w1}, \eqref{0-au'0} and \eqref{0-u0}.
\end{proof}

\end{document}
